%\documentclass[12pt]{article}
%\input{../../../../../Headers/TestPackageHeader}
%
%%\usepackage{geometry}
%\usepackage{amsmath}
%\usepackage{amssymb}
%\usepackage{amsthm}
%\usepackage{amscd}
%\usepackage{amsfonts}
%\usepackage{graphicx}%
%\usepackage{array}
%\usepackage{pgf,tikz}
%\usetikzlibrary{arrows}
%
%% Uncomment the following lines to use the Palatino font.  Remove the
%% [osf] bit if you don't like the old style figures.
%
%%\usepackage[T1]{fontenc}
%%\usepackage[osf]{mathpazo}
%
%\renewcommand{\headrulewidth}{0pt} % remove lines as well
%
%\numberwithin{equation}{section}
%
%\DeclareGraphicsRule{.tif}{png}{.png}{`convert #1 `dirname #1`/`basename #1 .tif`.png}
%
%\textheight = 10.8 in
%\topmargin = -1.0in
%\textwidth = 7.7in
%\oddsidemargin = -0.7in
%\evensidemargin = -0.7in
%\newcommand{\A}{{(\alpha)}}
%\newcommand{\B}{{(\beta)}}
%\newcommand{\K}{\kappa^{(i)}}
%
%\usepackage{sectsty}
%\subsectionfont{\rmfamily\mdseries\itshape\large}


\documentclass[11pt]{article}

% Package Header

\usepackage[paperheight=11in,
			paperwidth=8.5in,
			inner=0.4in,
			textheight=8.9in,
			textwidth=7.8in,
			marginparwidth=120pt,
			marginparsep=32pt,
			bottom=1.4in,
			footskip=1.0in]{geometry}
\geometry{letterpaper}                   % ... or a4paper or a5paper or ...

\usepackage{array}
\usepackage{amssymb}
\usepackage{asymptote}
\usepackage{amsmath}
\usepackage{amsfonts}
\usepackage{animate}
\usepackage{arrayjobx}
\usepackage{arydshln}
\usepackage{bookman}
\usepackage{color}
%\usepackage{colortbl}
\usepackage{cancel}
\usepackage{calculator}
\usepackage{collcell}
\usepackage{epstopdf}
\usepackage{enumerate}
\usepackage{epsfig}
%\usepackage{etex}
\usepackage[english]{babel}
\usepackage{etoolbox}
% \usepackage{fancyhdr}
\usepackage{fancybox}
\usepackage[nomessages]{fp}
\usepackage{graphics}
\usepackage{graphicx}
\usepackage{ifthen}
\usepackage{ifmtarg}
\usepackage[latin1]{inputenc}
\usepackage{latexsym}
\usepackage{multirow}
\usepackage{multicol}
\usepackage{multimedia}
\RequirePackage[noplaybutton]{media9}
\usepackage{mathbbol}
\usepackage{mathtools}
\usepackage{multido}
\usepackage{palatino}
\usepackage{polynom}
\usepackage{pxfonts}
\usepackage{pifont}
\usepackage{pgffor}
%\usepackage{pgfplots}
\usepackage{pdfpages}
\usepackage[parfill]{parskip}
\usepackage{soul}
\usepackage{tikz}
\usetikzlibrary{arrows,shapes,trees,shapes.geometric,calc}
\usepackage{txfonts}
\usepackage{tcolorbox}
\usepackage{wasysym}
\usepackage{wrapfig}
\usepackage{xcolor}
\usepackage{xmpmulti}
\usepackage{xparse}
\usepackage{xstring}
\usepackage{enumitem}

%\usepackage{hyperref}
\usepackage{geometry}
\usepackage{amsmath}
\usepackage{amssymb}
\usepackage{amsthm}
\usepackage{amscd}
\usepackage{amsfonts}
\usepackage{graphicx}%
\usepackage{fancyhdr}
\usepackage{array}
\usepackage{pgf,tikz}
\usetikzlibrary{arrows}

% Uncomment the following lines to use the Palatino font.  Remove the
% [osf] bit if you don't like the old style figures.

%\usepackage[T1]{fontenc}
\usepackage[osf]{mathpazo}
%\usepackage{mathpazo}

% Set your name here
%\def\name{Geoffrey W. Cox}

% The following metadata will show up in the PDF properties
%\hypersetup{
%  colorlinks = true,
%  urlcolor = black,
%  pdfauthor = {\name},
%  pdfkeywords = {economics empirical industrial organization econometrics
%    microeconomics applied microeconomics},
%  pdftitle = {\name: Curriculum Vitae},
%  pdfsubject = {Curriculum Vitae},
%  pdfpagemode = UseNone
%}

\geometry{textheight=9.4in, textwidth=7.7in}

%% Customize page headers
%\pagestyle{myheadings}
%\markright{\name}
%\thispagestyle{empty}

% Customize section headings
\usepackage{sectsty}
\subsectionfont{\rmfamily\mdseries\itshape\large}

% Don't indent paragraphs.
\setlength\parindent{0em}


\begin{document}

\input{../../../../../Headers/TestDefinitionHeader.tex}

$\quad$

\vspace{1.0cm}

\begin{center}
\LARGE {MA 301} \\

\Large{Test 1} \\

\large{Spring - 2013} \\

\vspace{0.5cm}

\setlength{\extrarowheight}{0.4cm}
\begin{tabular}{|c|c|c|}
\hline
Problem & points & out of  \\
\hline
\hline
1 & & 10 \\
\hline
2 & & 10 \\
\hline
3 & & 10 \\
\hline
4 & & 10 \\
\hline
5 & & 10 \\
\hline
6 & & 10 \\
\hline
7 & & 10 \\
\hline
8 & & 10 \\
\hline
\hline
Score & & 80\\
\hline
\end{tabular}
\vspace{0.5cm}

%--------TABLE---------
	\setlength{\extrarowheight}{0.1cm}
	\begin{tabular}{rcc}
		Name: & \underline{$\qquad \qquad \qquad \qquad \quad$}& \\
		&&$\hspace{1.4cm}$\\
		HR: & \underline{$\qquad \qquad \qquad \qquad \quad$}& \\
		&&\\
		Signature: & \underline{$\qquad \qquad \qquad \qquad \quad$}& \\
	\end{tabular}
%--------------------

\vspace{0.5cm}

	\setlength{\extrarowheight}{0.1cm}
	\begin{tabular}{| p{12cm} |}
		\hline
		\multicolumn{1}{|c|}{{\bf READ BEFORE YOU START}} \\
		\hline
		 * Show all work.\\
		 * Change your calculator mode to RADIANS.\\
		 * Sloppy solutions may be penalized.\\
		 * Use calculators ONLY for computations.\\
		 * You must write up every problem in a clear and \\
		 $\,\,$ organized way.\\
		 * If you write solutions on a page different than the \\
		 $\,\,$ question, then clearly state where it is located.\\
		\hline
	\end{tabular}

\end{center}

\vspace{1.0cm}

\thispagestyle{empty}
\newpage
\thispagestyle{empty}

{\bf Problem 1:} Compute \,\,$\D{\OWORK}$ \,\, if ${\bf F} = [x^2 + 4y, x + y^2]$ and $C$ is the rectangle in the $xy-$plane with $0 \le x \le 4$ and $-1 \le y \le 2$

{\bf Problem 2:} Let \, $\D{{\bf F} = \left[\frac{y^2}{x},\,2y\ln x,\,3z^2\right]}$ compute the line integral over $C$ shown below. Justify your work.
\begin{center}
	\includegraphics[scale = 0.38]{../../../Figures/Test1_Fig3.png}
\end{center}

{\bf Problem 3:} Let ${\bf F} = [-y,x]$ and $C$ be a simple closed curve in the $xy-$plane. Let $R$ be the region inside $C$.
 
(a) Show that $\D{\text{Area}(R) = \frac{1}{2}\oint \limits_{C} \text{\bf F} \cdot d\text{\bf r}}$

(b) Use part (a) to find the area inside the ellipse $C$: $4x^2 + y^2 = 4$ \,\,({\em Do NOT use a double integral to get the area})

\vspace{0.2cm}

{\bf Problem 4:}  \textcolor{blue}{\bf Set-up} the integral that gives the flux of water through the parabolic cylinder $S:\, y=x^2$ where $0 \le x \le 2$ and $0 \le z \le 3$ if the velocity flow of water is ${\bf v} = [3z^2, 6, 6xz]$ (units: {\em m/s}).  (water density: $\rho = 1 \, kg/m^3$).

\vspace{0.1cm}

%Compute the surface integral \qquad $\D{\iint \limits_{S} (1 + 9xz)^{3/2}\, dA}$
%
%\vspace{0.1cm}
%
%where $S$ is given by the parameterization: ${\bf r}(u,v) = [u,v,u^3]$,\,\, $0 \le u \le 1$,\,\, $-2 \le v \le 2$
%
%\vspace{0.2cm}
%
%Compute the surface integral of ${\bf F} = [e^y, -e^z, e^x]$ over the surface $S:\, x^2 + y^2 = 25$ \, for $x \ge 0$, $y \ge 0$, and $0 \le z \le 2$
%
%\vspace{0.2cm}

{\bf Problem 5:}  The closed circular cylinder shown below has a height of 10 and a base radius of 6. Let $A$ be the top, $B$ the bottom, and $S$ the side surface of the cylinder. Also, let ${\bf F} = (7x+1){\bf i} - 2y{\bf j} + z{\bf k}$. Find the outward flux through side $S$.  ({\em Note}:  $Vol \, of \, a \, cylinder = \pi r^2 h$ and $Area \, of \, a \, circle = \pi r^2$)
\begin{center}
	\includegraphics[scale = 0.28]{../../../Figures/Test1_Fig1.png}
\end{center}

{\bf Problem 6:} Let $\text{\bf F} = [x,y,z]$.   Compute the flux of $\text{\bf F}$ through the following surfaces $S$:

(a) $S$ is the hemi-sphere $x^2+y^2+z^2 = 25$, $z \ge 0$ with outward orientation. (\textcolor{blue}{\bf Set-up Only})

\vspace{0.1cm}

(b) $S$ is the surface of the rectangular box: $|x| \le 1$, \,\, $|y| \le 3$, \,\, and $0 \le z \le 2$. (\textcolor{blue}{\bf Compute this one})

\vspace{0.2cm}

{\bf Problem 7:} \textcolor{blue}{\bf Set-up} the integrals needed to compute the following:

(a) The work done by $\text{\bf F} = [xy,2y^2]$ in moving a mass along $\text{\bf r}(t) = [\cos t, \sin t]$ \, \& \, $\D{0 \le t \le \frac{\pi}{2}}$.
	
(b) The mass of a wire parameterized by \RT = t\VI + 2t\VJ + $\frac{2}{3}t^{3/2}$\VK \,where\, $0 \le t \le 2$ with density $\delta(x,y,z) = 3\sqrt{5+x}$.

\vspace{0.01cm}

(c) The surface area of $S$ with parameterization: ${\bf r}(u,v) = [u,v,u^3]$,\,\, $0 \le u \le 1$,\,\, $-2 \le v \le 2$.

{\bf Problem 8:} Answer the following true/false questions

\vspace{0.1cm}

	\setlength{\extrarowheight}{0.2cm}
	\begin{tabular}{| c | p{17.2cm} |}
		\hline
		T\quad/\quad F & The line integral of any vector field {\bf F} around any closed curve $C$ is zero\\
		\hline
		T\quad/\quad F & If $\WORK = 0$ for one particular closed path, the ${\bf F}$ is path-independent.\\
		\hline
		T\quad/\quad F & If ${\bf F}$ is path-independent, then there is a potential function for ${\bf F}$\\
		\hline
		T\quad/\quad F & If $\WORK \ne 0$ for some closed path $C$, then ${\bf F}$ is path-dependent.\\
		\hline
		T\quad/\quad F & If ${\bf F} = \nabla f$ then ${\bf F}$ is conservative.\\
		\hline
		T\quad/\quad F & If ${\bf F}$ is path-independent and $C$ is any closed curve, then $\WORK = 0$\\
		\hline
%		T\quad/\quad F & If $S_1$ is a rectangle with area = 1 and $S_2$ is a rectangle with area = 2, then
%$2\iint \limits_{S_1} {\bf F} \cdot d{\bf A} = \iint \limits_{S_2} {\bf F} \cdot d{\bf A}$\\
%		\hline
		T\quad/\quad F & Let $S$ be a sphere of radius 1 centered at the origin and oriented outward.  If ${\bf F} = [x,y,z]$ then $\FLUX$ is positive\\
		\hline
		T\quad/\quad F & Let $S$ be a sphere of radius 1 centered at the origin and oriented outward.  If $\FLUX = 0$, then ${\bf F} = {\bf 0}$\\
		\hline
%		T\quad/\quad F & Let $S$ be the surface of a cube centered at the origin with side length 2 and oriented outward.  If ${\bf F} = 2\VI - \VK$ then the flux of ${\bf F}$ through $S$ is 27.\\
%		\hline
		T\quad/\quad F & If ${\bf F}$ and ${\bf G}$ are vector-fields satisfying div(${\bf F}$) = div(${\bf G}$), then ${\bf F} = {\bf G}$\\
		\hline
		T\quad/\quad F & If ${\bf F}$ is a vector-field such that div(${\bf F}$) = 1 and $S$ is a closed surface oriented outward, then $\FLUX$ is equal to the volume enclosed by $S$.\\
		\hline
%		T\quad/\quad F & If $\FLUX = 12$ and $S$ is a flat disk of area $4\pi$, then div$({\bf F}) = \frac{3}{\pi}$\\
%		\hline
	\end{tabular}

\thispagestyle{empty}
\newpage
\thispagestyle{empty}

\end{document}